\section{CloudTrail}
\begin{itemize}
    \item Provides governance, compliance and audit for your AWS Account
    \item CloudTrail is enabled by default.
    \item Get a history of events / API calls made within your AWS Account by;
    \item Console
    \item SDK
    \item CLI
    \item AWS Services
    \item Can put logs from CloudTrail into CloudWatch Logs or S3.
    \item A trail can be applied to All Regions (default) or a single Region.
    \item If a resource is deleted in AWS, investigate CloudTrail first
\end{itemize}

%% TODO Insert cloud trail diagram here

\begin{itemize}
    \item Management Events:
    \item Operations that are performed on resources in your AWS account
    \begin{itemize}
        \item Examples:
        \item Configuring security (IAM AttachRolePolicy)
        \item Configuring rules for routing data (Amazon EC2 CreateSubnet)
        \item Setting up logging (AWS CloudTrail CreateTrail)
        \item By default, trails are configured to log management events.
        \item Can separate Read Events (that don't modify resources) from Write Events (that may modify resources)
    \end{itemize}

    \item Data Events:
    \begin{itemize}
        \item By default, data events are not logged (because high volume operations)
        \item Amazon S3 object-level activity (ex: GetObject, DeleteObject, PutObject): can separate Read and Write Events
        \item AWS Lambda function execution activity (the Invoke API)
    \end{itemize}

    \item CloudTrail Insights
    \begin{itemize}
        \item Enable CloudTrail Insights to detect unusual activity in your account:
        \item Inaccurate resource provisioning
        \item Hitting service limits
        \item Burst of AWS IAM actions
        \item Gaps in periodic maintenance activity
        \item CloudTrail Insights analyses normal management events to create baseline
        \item And then continuously analyses write events to detect unusual patterns
        \item Anomalies appear in the CLoudTrail console
        \item Event is sent to Amazon S3
        \item An EventBridge event is generated (for automation needs)
    \end{itemize}

%TODO Diagram

    \item CloudTrail Events Retention
    \begin{itemize}
        \item Events are stored for 90 days in CloudTrail
        \item To Keep events beyond this period, log them to S3 and use Athena
    \end{itemize}
\end{itemize}

%TODO Diagram


\section{CloudTrail - EventBridge Integration}

%TODO Diagram Amazon EventBridge - Intercept API calls
%TODO Diagram Amazon EventBridge - CloudTrail


\section{CloudTrail - SA Pro}
%Add diagram

\subsection{S3 Enhancements}:
\begin{itemize}
    \item Enable Versioning
    \item MFA Delete Protection
    \item S3 Lifecycle Policy (S3 IA, Glacier)
    \item S3 Object Lock
    \item SSE-S3 or SSE-KMS encryption
    \item Feature to perform CloudTrail Log File Integrity validation (SHA-256 for hashing and signing)
\end{itemize}

%TODO Diagram - CloudTrail - Solution Architecture - Multi Account, Multi Region Logging

\subsection{Observations}
\begin{itemize}
    \item The S3 Bucket policy is necessary for cross-account delivery
    \item If Account A wants to access its CloudTrail files:
    \item Option 1: Create a cross-account role and assume the role
    \item Option 2: edit the bucket policy
\end{itemize}

%TODO - CloudTrail - Solution Architecture - Alert for API calls
- Log filter metrics can be used to detect a high level of API happening
- Ex: Count occurrences of EC2 TerminateInstances API
- Ex: Count of API calls per user
- Ex: Detect high level of Denied API calls

%TODO - Diagram - CloudTrail - Solution Architecture: Organisational Trail
- The Organisational Trail is created in the management account.

\subsection{CloudTrail: How to react to events the fastest?}
\begin{itemize}
    \item Overall, CloudTrail may take up to 15 minutes to deliver events
    \item EventBridge:
    \item Can be triggered for any API call in CloudTrail
    \item The fastest, most reactive way.
    \item CloudTrail Delivery in CloudWatch Logs:
    \item Events are streamed
    \item Can perform a metric filter to analyse occurrences and detect anomalies
    \item CloudTrail Delivery in S3
    \item Events are delivered every 5 minutes
    \item Possibility of analysing logs integrity, deliver cross account, long-term storage.
\end{itemize}


\section{KMS}

\subsection{AWS KMS (Key Management Service)}
\begin{itemize}
    \item Anytime you hear "encryption" for an AWS service, it is most likely KMS
    \item Easy way to control access to your data, AWS manages keys for us.
    \item Fully integrated with IAM for authorisation
    \item Seamlessly integrated into:
    \item Amazon EBS: encrypt volumes
    \item Amazon S3: Server-side encryption of objects
    \item Amazon Redshift: Encryption of data
    \item Amazon RDS: Encryption of data
    - Etc.....
    \item Can also use the CLI / SDK
\end{itemize}

\subsection{KMS - KMS - Key Types}
\begin{itemize}
    \item Symmetric (AES-256 keys)
    \item First offering of KMS, single encryption key that is used to Encrypt and Decrypt
    \item AWS services that are integrated with KMS use Symmetric KMS keys
    \item Necessary for envelope encryption
    \item You never get access to the KMS key unencrypted (must call KMS API to use)
    \item Asymmetric (RSA and ECC key pairs)
    \item Public (Encrypt) and Private Key (Decrypt) pair
    \item Used for Encrypt/Decrypt or Sign/Verify operations
    \item The public key is downloadable but you can't access the Private Key unencrypted
    \item Use case: encryption outside of AWS by users who can't call the KMS API
\end{itemize}

\subsection{Types of KMS Keys}
\begin{itemize}
    \item Customer Managed Keys
    \item Create, manage and use. Can enable or disable
    \item Possibility of rotation policy (new key generated every year, old key preserved)
    \item Can add a Key Policy (resource policy) and audit in CloudTrail
    \item Leverage for envelope encryption.
    \item AWS Managed Keys
    \item Used by AWS service (aws/s3, aws/ebs, aws/redshift)
    \item Managed by AWS (automatically rotated every 1 year)
    \item View Key Policy and audit in CloudTrail
    \item AWS Owned Keys
    \item Create and managed by AWS, use by some AWS services to protect your resources
    \item Used in multiple AWS accounts but they are no in your AWS account
    \item You can't view, use, track or audit.
\end{itemize}

Types of KMS keys
%TODO Insert table

\subsection{KMS Key Material Origin}
- Identifies the source of the key material in the KMS key
- Can't be changed after creation

- KMS(AWS_KMS) - default
- AWS KMS created and managed the key material in its own key store

- External (EXTERNAL)
- You import the key material into the KMS key
- You're responsible for securing and managing this key material outside of AWS

- Custom Key Store (AWS_CLOUDHSM)
- AWS KMS creates the key material in a custom key store (CloudHSM Cluster)

\subsection{KMS Key Source - Custom Key Store (CloudHSM)}
\begin{itemize}
    \item Integrate KMS with CloudHSM cluster as a Custom Key Store
    \item Key materials are stored in a CloudHSM cluster that you own and manage
    \item The cryptographic operations are performed in the HSMs
    \item Use cases:
    \item You need direct control over the HSMs
    \item KMS Keys needs to be stored in a dedicated HSMs
\end{itemize}

\subsection{KMS Key Store - External}
\begin{itemize}
    \item Import your own key material into KMS key, Bring Your Own Key (BYOK)
    \item You're responsible for key material's security, availability and durability outside of AWS
    \item Supports both Symmetric and Asymmetric KMS keys
    \item Can't be used with Custom Key Store (CloudHSM)
    \item Manually rotate your KMS key (Automatic and On-demand Key Rotation are NOT supported)
\end{itemize}

%TODO Diagram

\subsection{KMS Multi-Region Keys}
%TODO Diagram
\begin{itemize}
    \item A set of identical KMS keys in different AWS Regions that can be used interchangeably (same KMS key in multiple Regions)
    \item Encrypt in one Region and decrypt in other Regions (No need to re-encrypt or making cross-Region API calls)
    \item Multi-Region keys have the same key ID, key material, automatic rotation........
    \item KMS - Multi-Region are NOT global (Primary + Replicas)
    \item Each Multi-Region key is manged independently
    \item Only one primary key at a time, can promote replicas into their own primary
    \item Use cases: Disaster Recovery, Global Data Management (e.g DynamoDB, Global Tables), Active-Active Applications that span multiple Regions, Distributed Signing applicatoins.
\end{itemize}


\section{Parameter Store}

\subsection{SSM Parameter Store}
\begin{itemize}
    \item Secure storage for configuration and secrets
    \item Optional Seamless Encryption using KMS
    \item Serverless, scalable, durable easy SDK
    \item Version tacking of configurations / secrets
    \item Security through IAM
    \item Notifications with Amazon EventBridge
    \item Integration with CloudFormation
\end{itemize}

%TODO Insert diagram

SSM Parameter Store Hierarchy

%TODO Insert diagram

Standard and advanced parameter tiers

%TODO tables

Parameter Policies (for advanced parameters)
\begin{itemize}
    \item Allows to assign a TTL to a parameter (expiration date) to force updating or deleting sensitive data such as passwords.
    \item Can assign multiple policies at a time.
\end{itemize}


\section{Secrets Mangager}

\subsection{AWS Secrets Manager}

\begin{itemize}
    \item Meant for storing secrets (e.g passwords, API keys,......)
    \item Capability to force rotation of secrets every X days
    \item Automate generation of secrets on rotation (uses Lambda)
    \item Natively supports Amazon RDS (all supported DB engines)
    \item Support other databases and services (custom Lambda function)
    \item Control access to secrets using Resource-based Policy
    \item Integration with other AWS services to natively pull secrets from Secrets Manager: CloudFormation, CodeBuild, ECS, EMR, Fargate, EKS, Parameter Store......
\end{itemize}
%TODO Insert diagram

%TODO insert code listing ---- Secrets Manager with CloudFormation

%TODO insert diagram ---- Secrets Mangager - Sharing Across Accounts

\subsubsection{SSM Parameter Store vs Secrets Manager}
\begin{itemize}
    \item  Secrets Manager (dollah)
    \item  Automatic rotation of secrets with AWS lambda
    \item  Lambda function is provided for RDS, Redshift, DocumentDB
    \item  KMS encryption is mandatory
    \item  Can integrate with CloudFormation
\end{itemize}

\subsubsection{SSM Parameter Store (dollah)}
\begin{itemize}
    \item Simple API
    \item No secret rotation (can enable rotation using Lambda triggered by EventBridge)
    \item KMS encryption is optional
    \item Can integrate with CloudFormation
    \item Can pull a Secrets Manager secret using the SSM parameter Store API
\end{itemize}

SSM Parameter Store vs Secrets Manager Rotation
%TODO diagram


\section{RDS Security}
\begin{itemize}
    \item KMS encryption at rest for underlying EBS volumes / snapshots
    \item Transparent Data Encryption (TDE) for Oracle and SQL Server
    \item SSL Encryption to RDS is possible for all DB (in-flight)
    \item IAM Authentication for MySQL, PostgreSQL and MariaDB
    \item Authorisation still happens within RDS (not in IAM)
    \item Can copy an un-encrypted RDS snapshot into an encrypted one
    \item CloudTrail cannot be used to track queries made within RDS
\end{itemize}


\section{SSL Encryption, SNI, MITM}

\subsection{SSL/TLS - Basics}
\begin{itemize}
    \item SSL refers to Secure Sockets Layer, used to encrypt connections
    \item TLS refers to Transport Layer Security which is a newer version
    \item Nowadays, TLS certificates are mainly used but people still refer as SSL
    \item Public SSL certificates are issued by Certificate Authorities (CA)
    \item Comodo, Symantec, GoDaddy, GlobalSign, Digicert, Letencrypt, etc
    \item SSL certificates have an expiration date (you set) and must be renewed
\end{itemize}

\subsection{SSL Encryption - How it works}
%TODO insert sequence diagram
\begin{itemize}
    \item Asymmetric Encryption is expensive (SSL)
    \item Symmetric encryption is cheaper
    \item Asymmetric handshake is used to exchange a per-client random symmetric key.
    \item Possibility of client sending an SSL certificate as well (two-way certificate)
\end{itemize}

\subsection{SSL - Server Name Indication (SNI)}
%TODO insert diagram
\begin{itemize}
    \item SNI solves the problem of loading multiple SSL certificates onto one web server (to server multiple websites)
    \item It's a "newer" protocol and requires the client to indicate the hostname of the target server in the inisital SSL handshake
    \item The server will the find the correct certificate or return the default one.
\end{itemize}

Note:
\begin{itemize}
    \item Only works for ALB and NLB (newer generation), CloudFront
    \item Does not work for CLB (older gen)
\end{itemize}

\subsection{SSL - Man in the middle attacks}
%TODO - add sequence diagram
How to prevent
\begin{enumerate}
    \item Don't use public-facing HTTP, use HTTPS (meaning use SSL/TLS certificates)
    \item Use a DNS that has DNSSEC
    \begin{itemize}
        \item To send a client to a pirate server, a DNS response needs to be "forged" by a server which intercepts them.
        \item It is possible to protect your domain name by configuring DNSSEC
        \item Amazon Route 53 supports DNSSEC for domain registration
        \item Route 53 supports DNSSEC for DNS service as of December 2020 (using KMS)
        \item You could also run a custom DNS server on Amazon EC2 for example (Bind is the most popular, dnsmasq, KnotDNS, PowerDNS)
    \end{itemize}
\end{enumerate}


\section{AWS Certificate Manager - ACM}

\subsection{AWS Certificate Manager (ACM)}
\begin{itemize}
    \item To host public SSL certificates in AWS, you can
    \item Buy your own and upload them using the CLI
    \item Have ACM provision and renew public SSL certificates for you (free of cost)

    \item ACM loads SSL certificates on the following integrations:
    \item Load Balancers (including the ones created by EB)
    \item CloudFront distributions
    \item APIs on API gateways

    \item SSL certificates is overall a pain to manually manage, so ACM is great to leverage in your AWS infrastructure
\end{itemize}

%TODO Diagram

\subsection{ACM - Good to know}
\item Possibility of creating public certificates
\item Must verify public DNS
\item Must be issued by a trusted public certificate authority (CA)
\item Possibility of creating private certificates
\item For you internal applications
\item You create your own private CA
\item Your applications must trust your private CA
\item Certificate renewal:
\item Automatically done if generated provisioned by ACM
\item Any manually uploaded certificates must be renewed manually and re-uploaded.
\item ACM is a regional service
\item To use with a global application (multilpe ALB for example), you need to issue an SSL certificate
\item You cannot copy certs across regions


\section{CloudHSM}
\begin{itemize}
    \item KMS => AWS manages the software for encryption
    \item CloudHSM -> AWS provisions encryption hardware
    \item Dedicated Hardware (HSM = Hardware Security Module)
    \item You manage your own encryption keys entirely (not AWS)
    \item HSM device is tamper resistant, FIPS 140-2 Level 3 compliance
    \item Supports both symmetric and asymmetric encryption (SSL/TLS keys)
    \item No free tier available
    \item Must use the CloudHSM Client Software
    \item Redshift supports CloudHSM for database encryption and key management
    \item Good option to use with SSE-C encryption
\end{itemize}

\subsection{CloudHSM Diagram}
%TODO diagram
\begin{itemize}
    \item IAM permissions
    \item CRUD an HSM Cluster
    \item CloudHSM Software
    \item Manage the Keys
    \item Manage the Users
\end{itemize}

\subsection{CloudHSM - High Availability}
\item CloudHSM clusters are spread across Multi AZ (HA)
\item Great for availability

%TODO diagram

\subsection{CloudHSM vs KMS}
%TODO table


\section{Solution Architecture - SSL on ELB}
%TODO insert diagram

%TODO insert diagram

\begin{itemize}
    \item You can offload SSL to CloudHSM (SSL Acceleration)
    \item Supported by NGINX Apache Web servers and IIS for windows server
    \item Extra security: the SSL private key never leaves the HSM device
    \item Must setup a cryptographic user (CU) on the CloudHSM device
\end{itemize}

%TODO insert diagram


\section{S3 Security}
\begin{itemize}
    \item SSE-S3: encrypts S3 objects using keys handled and managed by AWS
    \item SSE-KMS: leverage KMS to manage encryption keys
    \item Key usage appears in CloudTrail
    \item objects made public can never be read
    \item On s3:PutObject make the permission kms:GenerateDataKey is allowed
    \item SSE-C: when you want to manage your own encryption keys
    \item Client-Side Encryption
\end{itemize}

Glacier: all data is AES-256 encrypted, key under AWS control

\subsection{Encryption in transit (SSL/TLS)}
- Amazon S3 exposes:
\begin{itemize}
    \item HTTP endpoint: non encrypted
    \item HTTPS endpoint: encryption
    \item You're free to use the endpoint you want, but HTTPS is recommended
    \item HTTPS is mandatory for SSE-C
    \item To enforce HTTPS, use a bucket policy wit aws:SecureTransport
\end{itemize}

\subsection{Events in S3 Buckets}
\begin{itemize}
    \item S3 Access Logs:
    \begin{itemize}
        \item Detailed records for the requests that are made to a bucket
        \item Might take hours to deliver
        \item Might be incomplete (best effort)
    \end{itemize}
    \item S3 Event Notifications
    \begin{itemize}
        \item Receive notifications when certain events happen in your bucket
        \item E.g: new objects created, object removal, restore objects, replication events.
        \item Destinations: SNS, SQS queue, Lambda
        \item Typically delivered in seconds but can take minutes, notification for every object if versioning is enabled, else risk of one notification for two same objects write done simultaneously.
    \end{itemize}
    \item Trusted Advisor:
    \begin{itemize}
        \item Check the bucket permission (is the bucket public?)
        \item Amazon EventBridge
        \item Need to enable CloudTrail object level logging on S3 first
        \item Target can be Lambda, SQS, SNS, etc.........
    \end{itemize}
\end{itemize}

\subsection{S3 Security}
\begin{itemize}
    \item User based
    \begin{itemize}
        \item IAM policies - which API calls should be allowed for a specific user from IAM console
    \end{itemize}

    \item Resource Based
    \begin{itemize}
        \item Bucket Policies - bucket wide rules from the S3 console - allows cross account
        \item Object Access Control List (ACL) - finer grain
        \item Bucket Access Control List (ACL) - less common
    \end{itemize}
\end{itemize}

\subsection{S3 Bucket Policies}
\begin{itemize}
    \item Use S3 bucket for policy to:
    \item Grant public access to the bucket
    \item Force objects to be encrypted at upload
    \item Grant access to another account (Cross Account)
    \item Optional Conditions on:
    \item SourceIp: Public IP or Elastic IP | VpcSourceIp: Private IP (through VPC Endpoint)
    \item Source VPC or Source VPC Endpoint - only workds with VPC Endpoints
    \item CloudFront Origin Identity
    \item MFA
\end{itemize}
%TODO Examples --- insert link

\subsection{S3 pre-signed URLs}
\begin{itemize}
    \item Can generate pre-signed URLs using SDK or CLI
    \item For downloads (easy can use the CLI)
    \item For uploads (harder must use the SDK)
    \item Valid for a default of 3600 seconds, can change timeout with --expires-in[TIME_BY_SECONDS] argument
    \item Users given a pre-signed URL inherit the permissions of the person who generated the URL for GET / PUT
    \item Examples:
    \begin{itemize}
        \item Allow only logged-in users to download a premium video on your S3 bucket
        \item Allow an ever changing list of users to download files by generating URLs dynamically
        \item Allow temporarily a user to upload a file to a precise location in our bucket
    \end{itemize}
\end{itemize}

\subsection{VPC Endpoint Gateway for S3}
%TODO insert diagram

\subsection{S3 Object Lock and Glacier Vault Lock}
\begin{itemize}
    \item S3 Object Lock
    \item Adopt a WORM (Write once read many) model
    \item Block an object version deletion for a specified amount of time
    \item Glacier Vault Lock
    \item Adopt a WORM (Write Once Read Many) model
    \item Lock the policy for future edits (can no longer be changed)
    \item Helpful for compliance and data retention
\end{itemize}


\section{S3 Access Points}
%TODO insert diagram
\begin{itemize}
    \item Access Points simplify security management for S3 buckets
    \item Each access point has:
    \item its own DNS name (Internet Origin or VPC origin)
    \item an access point policy (similar to bucket policy) - manage security at scale
\end{itemize}

\subsection{S3 - Access Points - VPC Origin}
%TODO insert diagram
\begin{itemize}
    \item We can define the access point ot be accessible only form within the VPC
    \item You must create a VPC Endpoint to access the Access Point (Gateway or Interface Endpoint)
    \item The VPC Endpoint Policy must allow access to the target bucket and Access Point
\end{itemize}


\section{S3 Multi-Region Access Points}
%TODO insert a diagram
\begin{itemize}
    \item Provide a global endpoint that spans S3 buckets in multiple AWS regions
    \item Dynamically route requests to the nearest S3 bucket (lowest latency)
    \item Bi-directional S3 bucket replication rules are created to keep data in sync across regions
    \item Failover Controls - allows you to shift request across S3 buckets in different AWS regions within minutes (Active-Active or Active-Passive)
\end{itemize}

%TODO insert a diagram

Multi-Region Access Points - Failover Controls
%TODO insert diagram


\section{S3 Multi-Region Access Points - Hands On}
%TODO ------ or do via TerraForm :D


\section{S3 Object Lambda}
\begin{itemize}
    \item Use AWS Lambda Functions to change the object before it is retrieved by the caller application
    \item Only one S3 bucket is needed on top of which we create S3 Access Point and S3 Object Lambda Access Points
    \item Use Cases:
    \item Redacting personally identifiable information for analytics or non production environments
    \item Converting across data formats, such as converting XML to JSON
    \item Resizing and watermarking images on the fly using caller-specific details, such as the user who request the object
\end{itemize}

%TODO insert diagram


\section{DDoS and AWS Shield}
What is a DDoS
%TODO insert a diagram

\subsection{Types of Attacks on your infrastructure}
\begin{itemize}
    \item Distributed Denial os Service (DDoS)
    \item When your service is unavailable because it is receiving too many requests
    \item SYN Flood (Layer 4) - send too many TCP connection requests
    \item UDP reflection (Layer 4) - get other servers to send many big UDP requests
    \item DNS flood attack: Overwhelm the DNS so legitimate users can't find the site
    \item Slow Loris attack: a lot of HTTP connections are opened and maintained (hmmmm consider websocket servers!!)
\end{itemize}

\subsection{Application Level Attacks:}
\begin{itemize}
    \item more complex, more specific (HTTP level)
    \item Cache bursting strategies: overload the backend database by invalidating the cache (clever)
\end{itemize}

\subsection{DDoS Protection on AWS}
\begin{itemize}
    \item AWS Shield Standard: Protects against DDoS attack for your website and applications, for all customers at no additional cost
    \item AWS Shield Advanced: 24/7 premium DDoS protection
    \item AWS WAF: Filter specific requests based on rules
    \item CloudFront and Route 53:
    \item Availability protection using global edge network
    \item Combined with AWS shield: provides DDoS attack mitigation at the edge
    \item Be ready to scale - leverage AWS Auto Scaling
    \item Separate static resources (S3 / CloudFront) form dynamic ones (EC2 / ALB)
    \item Read the white paper %TODO insert link to white paper
\end{itemize}

\subsection{Sample Reference Architecture}
%TODO diagram

\subsection{AWS Shield}
\begin{itemize}
    \item AWS Shield Standard
    \item Free service that is activated for every AWS customer
    \item Provides protection from attacks such SYN/UDP Floods, Reflection attacks and other layer 3/layer 4 attacks
    \item AWS Shield Advanced
    \item Optional DDoS mitigation service (3000 dollahs per month per organisation)
    \item Protect against more sophisticated attack on Amazon EC2, Elastic Load Balancing (ELB), Amazon CloudFront, AWS Global Accelerator, Route 53
    \item 24/7 access to AWS DDoS response team (DRP)
    \item Protect against higher fees during usage spikes due to DDoS
\end{itemize}


\section{AWS WAF - Web Application Firewall}
\begin{itemize}
    \item Protects your web applications from common web exploits (Layer 7)
    \item Deploy on Application Load Balancer (localised rules)
    \item Deploy on API Gateway (rules running at the regional or edge level)
    \item Deploy on CloudFront (rules on edge locations)
    \item Used in front of other solutions: CLB, EC2 Instances, custom origins, S3 websites
    \item Deploy on AppSync (protect your GraphQL APIs)
    \item WAF is not for DDoS protection
    \item Define Web ACL (Web Access Control List)
    \item Rules can include IP addresses, HTTP headers, HTTP body, or URI strings
    \item Protects from common attack - SQL injection and Cross-Site scripting (XSS)
    \item Size constraints, Geo Match
    \item Rate-based rules (to count occurrences of events)
    \item Rule Actions: Count | Allow | Block | CAPTCHA | Challenge
\end{itemize}

\subsection{AWS WAF - Managed Rules}
\begin{itemize}
    \item Library of over 190 managed rules
    \item Ready to use rules that are managed by AWS and AWS Marketplace Sellers

    \item Baseline Rule Groups - general protection from common threats
    - AWSManagedRulesCommonRuleSet, AWSManagedRulesAdminProtectionRuleSet,.....
    \item Use-case Specific Rule Groups - Protection for many AWS WAF use cases
    - AWSManagedRulesSQLLiRuleSet , AWSManagedRulesWindowsRulesSet, AWSMangedRulesPHPRuleSet, AWSManagedRulesWordPressRuleSet
    \item IP Reputation Rule Groups - block requests based on source (e.g. malicious IPs)
    - AWSManagedRulesAmazonIpReputationList, AWSManagedRulesAnonymousIpList
    \item Bot Control Managed Rule Group - block and manage requests from bots
    - AWSMangedRulesBotControlRuleSet
\end{itemize}

\subsection{WAF - Web ACL - Logging}
\begin{itemize}
    \item You can send your logs to an:
    \item Amazon CloudWatch Logs log group - 5 MB per second
    \item Amazon Simple Storage Service (Amazon S3) bucket - 5 minutes interval
    \item Amazon Kinesis Data Firehouse - limited by Firehose quotas
\end{itemize}

%TODO diagrams

Solution Architecture - Enhance CloudFront Origin
Security with AWS WAF and AWS Secrets Manager

%TODO diagram


\section{AWS Firewall Manager}
\begin{itemize}
    \item Manage rules in all accounts of an AWS organisation
    \item Security policy: common set of security rules
    \item WAF rules (Application Load Balancer, API Gateways, CloudFront)
    \item AWS Shield Advanced (ALB, CLB, NLB, Elastic IP, CloudFront)
    \item Security Groups for EC2, Application Load Balancer and ENI resources in VPC
    \item AWS Network Firewall (VPC Level)
    \item Amazon Route 53 Resolver DNS Firewall
    \item Policies are created at the region level
\end{itemize}

Rules are applied to new resources as they are created (good for compliance) across all and future accounts in your Organisation

\subsection{WAF vs Firewall Manager vs Shield}
%Diagram
\begin{itemize}
    \item WAF, Shield and Firewall Manager are used together for comprehensive protection
    \item Define your Web ACL rules in WAF
    \item For granular protection of your resources, WAF along is the correct choice
    \item If you want to use AWS WAF across accounts, accelerate WAF configuration, automate the protection of new resources use Firewall Manager with AWS WAF
    \item Shield Advanced adds additional features on top of AWS WAF, such as dedicated support from the shield response team (SRT) and advanved reporting
    \item If you're prone to frequent DDoS attacks, consider purchasing Shield Advanved
\end{itemize}


\section{Blocking an IP Address}

\subsection{Blocking an IP address}
%TODO diagram

\subsection{Blocking an IP address - with an ALB}
%TODO diagram

\subsection{Blocking an IP address - with an NLB}
%TODO diagram

\subsection{Blocking an IP address - ALB + WAF}
%TODO diagram

\subsection{Blocking an IP address - ALB, CloudFront and WAF}
%TODO diagram


\section{Amazon Inspector}

\subsection{Amazon Inspector}
\begin{itemize}
    \item Automated Security Assessments
    \item For EC2 instances
    \item Leveraging the AWS System Manager (SSM) agent
    \item Analyse against unintended network accessibility
    \item Analyse the running OS against known vulnerabilities
    \item For container images push to amazon ECR
    \item Assessment of Container Images as they are pushed
    \item For Lambda Functions
    \item Identifies software vulnerabilities in function code and package dependencies
    \item Assessment of functions as they are deployed
    \item Reporting and integration with AWS security hub
    \item Send findings to Amazon Event Bridge
\end{itemize}

\subsection{What does Amazon Inspector evaluate?}
\begin{itemize}
    \item Remember: only for EC2 instances, Container Images and Lambda functions
    \item Continuous scanning of the infrastructure, only when needed
    \item Package vulnerabilities (EC2, ECR and Lambda) - database of CVE
    \item Network reachability (EC2)
    \item A risk score is associated with all vulnerabilities for prioritisation
\end{itemize}


\section{AWS Config}

\subsection{AWS Config}
\begin{itemize}
    \item Help with auditing and recording compliance of your AWS resources
    \item Helps records configuration and changes over time
    \item AWS Config rules does not prevent actions from happening (no deny)
    \item Questions that can be solved by AWS config
    \item Is there unrestricted SSH access to my security groups?
    \item Do my buckets have any public access?
    \item How has my ALB configuration changed over time?
    \item You can receive alerts (SNS notifications) for any changes
    \item AWS Config is a per-region service
    \item Can be aggregated across regions and accounts
\end{itemize}

\subsection{AWS config resource}
- View compliance of a resource over time
%TODO diagram

- View configuration of a resource over time
%TODO diagram

- View CloudTrail API calls if enabled
%TODO diagram

\subsection{AWS Config Rules}
\begin{itemize}
    \item Can use AWS managed config rules (over 75)
    \item Can make custom config rules (must be defined in AWS Lambda)
    \item Evaluate if each EBS disk is of type gp2
    \item Evaluate if each EC2 instance is t2.micro
    \item Rules can be evaluated / triggered
    \item For each config change
    And / or: at regular time intervals
    \item Trigger Amazon EventBridge if the rule is non-compliant (chain with Lambda)
    \item Rules can have auto remediation's through SSM Automations
    \item If a resource is not compliant you can trigger an auto remediation
    Ex: remediate security group rules, stop instances with non-approved tags
\end{itemize}


\section{AWS Managed Logs}
\begin{itemize}
    \item Load Balancer Access Logs (ALB, NLB, CLB) -> S3
    \item Access logs for your Load Balancers
    \item CloudTrail Logs -> to S3 and CloudWatch logs
    \item Logs for API calls made within your account
    \item VPC Flow Logs -> to S3, CloudWatch Logs, Kinesis Data Firehose
    \item Information about IP traffic going to and from network interfaces in your VPC
    \item Route 53 Access Logs -> to CloudWatch logs
    \item Log information about the queries that Route 53 receives
    \item S3 Access Logs -> to S3
    \item Server access logging provides detailed records for the requests that are made to a bucket
    \item CloudFront Access Logs -> to S3
    \item Detailed information about every user request that CloudFront receives
    \item AWS Config -> to S3
\end{itemize}


\section{Amazon GuardDuty}
\begin{itemize}
    \item Intelligent Threat discovery to protect your AWS Account
    \item Using Machine Learning algorithms, anomaly detection, 3rd party data
    \item One click to enable (30 days trial) no need to install software
    \item Input data includes:
    \item CloudTrail Events Logs - unusual API calls, unauthorised deployments
    \item CloudTrail Management Events - create VPC subnet, create trail
    \item CloudTrail S3 Data Events - get object, list objects, delete object
    \item VPC Flow logs - unusual internal traffic, unusual IP address
    \item DNS Logs - compromised EC2 instances sending encoded data within DNS queries
    \item Optional Feature - EKS Audit Logs, RDS and Aurora, EBS, Lambda, S3 Data Events.......
    \item Can setup EventBridge rules to be notified in case of findings
    \item EventBridge rules can target AWS Lambda or SNS
    \item Can protect against CryptoCurrency attacks (has a dedicated "finding" for it)
\end{itemize}

\subsection{Amazon GuardDuty}
%TODO diagrams

\subsection{GuardDuty - Delegated Administrator}
\begin{itemize}
    \item AWS organisation member accounts can be designated to be a GuardDuty Delegated Administrator
    \item Have full permissions to enable and manage GuardDuty for all accounts in the Organisation
    \item Can be done only using the Organisation Mangement Account
\end{itemize}
%TODO diagrams


\section{IAM Advanced Policies}

\subsection{IAM conditions}
%TODO - example JSON listing

%TODO - example JSON listing

\subsection{IAM for S3}
%TODO - example JSON listing
- s3:ListBucket permission applies to arn:aws:s3:::test -> bucket level permission

%TODO - example JSON listing
- s3:GetObject, s3:PutObject,s3:DeleteObject applies to arn:awn:s3:::test
- Object level permissions

\subsection{Resource Policies and aws:PrincipalOrgID}
- aws:PrincipalOrgID can be used in any resource polices to restrict access to accounts that are memeber of an AWS organisation
%TODO diagrram
%TODO json listing


\section{EC2 Instance Connect}
EC2 Instance Connect (SendSSHPublicKey API)
%TODO Diagram


\section{AWS Security Hub}
\begin{itemize}
    \item Central security tool to manage security across several AWS accounts and automate security checks
    \item Integrated dashboards showing current security and compliance status to quickly take actions
    \item Automatically aggregates alerts in predefined or personal findings formates from various AWS services and AWS partner tools:
    \item Config
    \item GuardDuty
    \item Inspector
    \item Macie
    \item IAM Access Analyser
    \item AWS Systems Manager
    \item AWS Firewall Manager
    \item AWS Health
    \item AWS Partner Network Solutions
\end{itemize}

Must first enable the AWS config service


\section{Amazon Detective}
\begin{itemize}
    \item GuardDuty, Macie and Security Hub are used to identify potential security issues or findings.
    \item Sometimes security findings require deeper analysis to isolate the root cause and take actions - it is a complex process
    \item Amazon Detective analyses, investigates and quickly identifies the root cause of security issues or suspicious activities (using ML and graphs)
    \item Automatically collects and processes events from VPC Flow Logs, CloudTrail, GuardDuty and creates a unified view.
    \item Produces visualisations with details and context to get to the root cause.
\end{itemize}
